\documentclass[12pt, a4paper]{article}
\usepackage[utf8]{inputenc}
\usepackage{amsmath, amssymb, amsthm, amsfonts}
\usepackage{CJKutf8}
\usepackage[margin=1in]{geometry}
\newtheorem{lemma}{Lemma}


\newenvironment{solution}
  {\paragraph{\textit{Solution:}}\par\medskip}
  {\hfill$\blacksquare$\vspace{0.1cm}}

\title{Topics in Geometry and Topology:Assignment 1}
\author{
    \begin{CJK*}{UTF8}{gbsn}
    钟皓宇 (12533556)
    \end{CJK*}
}
\date{\today}

\begin{document}

\maketitle
\section*{Problem 1}
If $S_{\sigma}$ is generated by $u_{1}, \dots, u_{t}$, and 
\[
A_{\sigma} = \mathbb{C}[\chi^{u_{1}}, \dots, \chi^{u_{t}}] = \mathbb{C}[Y_{1}, \dots, Y_{t}] / I,
\]
show that $I$ is generated by polynomials of the form
\[
Y_{1}^{a_{1}} Y_{2}^{a_{2}} \cdots Y_{t}^{a_{t}} - Y_{1}^{b_{1}} Y_{2}^{b_{2}} \cdots Y_{t}^{b_{t}}
\]
where $a_{1}, \dots, a_{t}, b_{1}, \dots, b_{t}$ are nonnegative integers satisfying the equation
\[
a_{1}u_{1} + \cdots + a_{t}u_{t} = b_{1}u_{1} + \cdots + b_{t}u_{t}.
\]
\begin{solution}
Let $\pi:\mathbb{C}[Y_1,\cdots,Y_t] \to A_{\sigma}, Y_{i} \mapsto \chi^{u_i}$, $I$ be the kernel of $\pi$ and $J$ be the ideal of $\mathbb{C}[Y_1,\cdots,Y_t]$ generated by the 
the form $Y_{1}^{a_{1}} Y_{2}^{a_{2}} \cdots Y_{t}^{a_{t}} - Y_{1}^{b_{1}} Y_{2}^{b_{2}} \cdots Y_{t}^{b_{t}}$,where $a_i,b_i\ge 0$ and $\sum a_iu_i = \sum b_iu_i$.

For an arbitrary generator $Y_{1}^{a_{1}} Y_{2}^{a_{2}} \cdots Y_{t}^{a_{t}} - Y_{1}^{b_{1}} Y_{2}^{b_{2}} \cdots Y_{t}^{b_{t}}$ of $J$, the ring homomorphism 
\begin{align*}
    \pi(Y_{1}^{a_{1}} Y_{2}^{a_{2}} \cdots Y_{t}^{a_{t}} - Y_{1}^{b_{1}} Y_{2}^{b_{2}} \cdots Y_{t}^{b_{t}}) = \chi^{\sum a_i u_i}- \chi^{\sum b_iu_i} = 0
\end{align*}
implies that $J\subseteq I$.

Suppose $f\in \mathbb{C}[Y_1,\dots,Y_t]$ be an arbitrary polynomial, then we have the graded decomposition of $f$:
\begin{align*}
    f =  \sum_{m\in S_\sigma} \sum_{\text{wt}(a)=m}c_a Y^{a}
\end{align*}
where $\text{wt}:\mathbb{Z}_{\ge 0}^*\to S_{\sigma},a=(a_1,\dots,a_t) \mapsto \sum a_iu_i$ is the weight function associated with $S_\sigma$ and $Y^a = Y_1^{a_1}\cdots Y_t^{a_t}$
is a monomial. If $f\in I$ is the polynomial in kernel of $\pi$, then we obtain that 
\begin{align*}
    \sum_{m\in S_\sigma}
    \left(\sum_{\text{wt}(a)=m}c_a \right)\chi^m = 0.
\end{align*} 
Interpret $\{\chi^m\}_{m\in S_\sigma}$ as a basis of vector space $\mathbb{C}[S_\sigma]$, then the coefficient 
\begin{align*}
    \sum_{\text{wt}(a)=m}c_a =0. \tag{$\star$}
\end{align*}
For any $m\in S_\sigma$, we select $a_{m}$ such that $\text{wt}(a_m)=m$. Using equation $(\star)$, we can rewrite the graded decomposition of $f$ by 
\begin{align*}
    f =  \sum_{m\in S_\sigma} \sum_{\text{wt}(a)=m}c_a (Y^{a}-Y^{a_m}).
\end{align*}
The graded decomposition we obtain implies $I\subseteq J$, our claim holds.
\end{solution}



\section*{Problem 2}
If $\sigma$ is a cone in $N$ and $\sigma'$ is a cone in $N'$, show that $\sigma \times \sigma'$ is a cone in $N \oplus N'$ and construct a canonical isomorphism
\[
U_{\sigma \times \sigma'} \cong U_{\sigma} \times U_{\sigma'}.
\]
\begin{solution}
Let $\{u_1,\dots,u_s\}$ and $\{v_1,\dots,v_t\}$ be the generators of the dual cones $\sigma^\vee$ and $(\sigma')^\vee$ of $\sigma$ and $\sigma'$ respectively.
Then, using the dual definition of cones, we obtain that
\begin{align*}
    \sigma \times \sigma' = \{w\in N_\mathbb{R}\oplus N_\mathbb{R}'|\langle w,(u_i,0)\rangle\ge 0\text{ for $i=1,\dots,s$}\text{ and }\langle w,(0,v_j)\rangle\ge 0\text{ for $j=1,\dots,t$}\},
\end{align*}
where we interpret the dual of $N\oplus N'$ by
\begin{align*}
    \operatorname{Hom}_{\mathbb{Z}}(N\oplus N',\mathbb{Z}) \cong \operatorname{Hom}_{\mathbb{Z}}(N,\mathbb{Z}) \oplus \operatorname{Hom}_{\mathbb{Z}}(N',\mathbb{Z}),
\end{align*}
and $\{(u_i, 0), (0, v_j)\}$ form a finite set of generators for the dual cone $(\sigma \times \sigma')^\vee$ inside $(N \oplus N')^\vee$. 
Therefore $\sigma \times \sigma'$ is expressed as the intersection of finitely many closed half-spaces and is, \textit{a fortiori}, a polyhedral cone.

The tensor product of the semigroup rings over $\mathbb{C}$ is isomorphic to the semigroup ring of their direct sum:
\begin{align*}
    A_\sigma \otimes_{\mathbb{C}} A_{\sigma'} = \mathbb{C}[S_\sigma] \otimes_{\mathbb{C}} \mathbb{C}[S_{\sigma'}] \cong \mathbb{C}[S_\sigma \oplus S_{\sigma'}] = A_{\sigma \times \sigma'}
\end{align*}
Taking the affine spectrum on both sides, we obtain the canonical isomorphism between the product of the affine toric varieties and the affine toric variety of the product cone:
$$U_\sigma \times U_{\sigma'} \cong \operatorname{Spec}(A_\sigma \otimes_{\mathbb{C}} A_{\sigma'}) \cong \operatorname{Spec}(\mathbb{C}[S_\sigma \oplus S_{\sigma'}]) = U_{\sigma \times \sigma'}.$$
Explicitly, the canonical $\mathbb{C}$-algebra isomorphism is given by extending the following map linearly:
\begin{align*}
    \varphi: A_\sigma \otimes_{\mathbb{C}} A_{\sigma'} &\xrightarrow{\sim} \mathbb{C}[S_\sigma \oplus S_{\sigma'}] = A_{\sigma \times \sigma'} \\
    \chi^u \otimes \chi^v &\mapsto \chi^{(u, v)}.
\end{align*}
Scheme-theoretically, we obtain a canonical isomorphism 
\begin{align*}
    \operatorname{Spec}(\varphi): U_{\sigma\times\sigma'} &\to U_{\sigma} \times U_{\sigma'} \\
    \mathfrak{p} &\mapsto \varphi^{-1}(\mathfrak{p}).
\end{align*}
\end{solution}



\section*{Problem 3}
\begin{enumerate}
    \item[(a)] Let $S \subset \mathbb{Z}_{\ge 0}$ be the sub-semigroup generated by 2 and 3. Then
    \[
    \mathbb{C}[S] = \mathbb{C}[X^{2}, X^{3}] = \mathbb{C}[Y, Z] / (Z^{2} - Y^{3}),
    \]
    so $\text{Spec}(\mathbb{C}[S])$ is a rational curve with a cusp.
    
    \item[(b)] Find $\mathbb{C}[S]$ when $S$ is the sub-semigroup of $\mathbb{Z}_{\ge 0}$ generated by a pair of relatively prime positive integers.
    
    \item[(c)] Find corresponding generators and relations for $S$ generated by 3, 5, and 7.
\end{enumerate}

\begin{solution}
\textbf{(a)} Let $\mathcal{X}=\operatorname{Spec}(\mathbb{C}[S])$. $\mathbb{C}[S]$ is an integral domain, because $Z^2-Y^3$ is irreducible over $\mathbb{C}[Y,Z]$, and zero ideal $(0)$ is a prime ideal, hence is a generic point of $\mathcal{X}$.
The function field of $\mathcal{X}$, as the stalk at generic point $(0)$, is $\operatorname{Frac}(\mathbb{C}[S])\cong \mathbb{C}(X)$. Therefore, 
by the transcendental degree of the function field, $\mathcal{X}$ is a curve, hence is rational because the function field of $\mathcal{X}$ is isomorphic to $\mathbb{C}(X)$.
By the Jacobian criterion of affine schemes, the Jacobian matrix 
\begin{align*}
    J = \begin{pmatrix} \frac{\partial f}{\partial Y} & \frac{\partial f}{\partial Z} \end{pmatrix} = \begin{pmatrix} -3Y^2 & 2Z \end{pmatrix}
\end{align*}
implies the only singular point of $\mathcal{X}$ is $(0,0)$. Consider the tangent cone of the singular point
\begin{align*}
    C_{(0,0)}(\mathcal{X}) = \operatorname{Spec}(\mathbb{C}[Y, Z] / (Z^2))
\end{align*}
is a non-reduced divisor of multiplicity 2 on $\mathbb{A}_{\mathbb{C}}^2$. Consequently, $\mathcal{X}$ is a rational curve with a cusp at $(0,0)$.

\textbf{(b)}
Suppose $(p,q)$ is an arbitrary coprime pair of positive integers. Consider a surjecteive ring homomorphism
\begin{align*}
    \pi: \mathbb{C}[Y,Z] \to \mathbb{C}[S], \quad
    Y\mapsto X^p, \quad
    Z\mapsto X^q
\end{align*}
Let $J=(Y^q-Z^p)$ be an ideal of $\mathbb{C}[Y,Z]$, then $\pi(Y^q-Z^p) = 0$ implies $J\subseteq \ker(\pi)$.
Since $\operatorname{gcd}(p,q)=1$, the polynomial $Y^q-Z^p$ is irreducible and therefore $\operatorname{Spec}(\mathbb{C}[Y,Z]/J)$ is an irreducible variety of dimension 1.
On the other hand, $\mathbb{C}[S]$ as the subring of $\mathbb{C}[X]$ is an integral domain, whence $\operatorname{Spec}(\mathbb{C}[S])$ is irreducible.
The function field of $\mathbb{C}[S]$, as the stalk at generic point $(0)$, is $\operatorname{Frac}(\mathbb{C}[S])\cong \mathbb{C}(X)$. 
By the transcendental degree of the function field of $\mathbb{C}[S]$, the dimension $\mathbb{C}[S]$ is 1. The inclusion relation of ideal $J\subseteq\ker(\pi)$ induces the inclusion relation of 
irreducible varieties of dimension 1, that is
\begin{align*}
    \operatorname{Spec}(\mathbb{C}[S]) \cong \operatorname{Spec}(\mathbb{C}[Y,Z]/\ker(\pi)) \subseteq \operatorname{Spec}(\mathbb{C}[Y,Z]/J).
\end{align*}
The irreducibililty and dimension force $\operatorname{Spec}(\mathbb{C}[S]) = \operatorname{Spec}(\mathbb{C}[Y,Z]/J)$.

\textbf{(c)}
Let $R = \mathbb{C}[Y, Z, W]$ and 
$
    \pi: R \to \mathbb{C}[S],
    Y\mapsto X^3,
    Z\mapsto X^5,
    W\mapsto X^7
$.
Let $I=\ker(\pi)$ and $f_1 = YW - Z^2,f_2 = Y^4 - ZW,f_3 = ZY^3 - W^2$. Let $J=(f_1,f_2,f_3)$, then $\pi(f_i)=0$ implies $J\subseteq I$. 
For any polynomial $g\in R$, we claim that $g\equiv\sum_{i}a_iY^i+\sum_{j}b_jY^jZ+\sum_{k}c_kY^kW\mod J$. Equivalently, for any element $\bar{g}\in R/J$ can be 
represented by a linear combination of $\{Y^i,Y^jZ,Y^kW\}$. Explicitly, we consider a monomial $Y^rZ^sW^t$. From relation $YW\equiv Z^2\mod J$, we can 
reduce the monomial to two cases: $Y^rZ^s$ or $Y^rW^t$. The first one can be reduced to $Y^rZ$ with remainer terms in the second case by $Z^2\equiv YW\mod J$. 
For $t\ge 2$, we can reduce $Y^rW^t$ to $Y^rW$ with remainer terms $Y^rZ^s$. This reduction can be terminates in finite steps, because the degree of $W$ is strictly decressing.
Consider canonical ring homomorphism $\gamma:R/J\to R/I$, suppose $\bar{g} = \sum_{i}a_iY^i+\sum_{j}b_jY^jZ+\sum_{k}c_kY^kW\mod J\in \ker(\gamma)$. Then 
\begin{align*}
    \gamma(\bar{g}) = \sum a_i X^{3i} + \sum b_jX^{3j+5}+\sum c_kX^{3k+7} = 0
\end{align*}
implies all coefficient $a_i,b_j,c_k$ is zero, because the $\{3i\},\{3j+5\},\{3k+7\}$ are pairwise disjoint, whence $\ker(\gamma)=I/J=(0)$ and $I=J$.
Therefore, we obtain that 
\begin{align*}
    \mathbb{C}[S] \cong  \mathbb{C}[Y, Z, W]/(YW - Z^2,Y^4 - ZW,ZY^3 - W^2).
\end{align*} 
\end{solution}



\section*{Problem 4}
If $\sigma \subset \mathbb{R}^{2}$ is the cone generated by $e_{2}$ and $\lambda \cdot e_{1} - e_{2}$, with $\lambda$ an irrational positive number, show that $\sigma^{\vee} \cap \mathbb{Z}^{2}$ is not finitely generated, and that $\mathbb{C}[\sigma^{\vee} \cap \mathbb{Z}^{2}]$ is not Noetherian.

\begin{solution}
Suppose $S = \sigma^{\vee}\cap \mathbb{Z}^2$ are finitely generated by $\{s_1,\dots,s_n\}$. By the definition of dual cones, we obtain that
\begin{align*}
    \sigma^\vee = \{(u_1,u_2)\in \mathbb{R}^2\mid  \lambda u_1\ge u_2\ge 0 \}.
\end{align*} 
For every point $u=(u_1,u_2)\in S$, we define a slope function $\Delta(u) = u_2/u_1$ and distence function
\begin{align*}
    D:S\to \mathbb{R}\quad 
    (u_1,u_2) \mapsto \frac{|\lambda u_1-u_2|}{\sqrt{1+\lambda^2}}.
\end{align*}
which evaluates the Euclidian distence of point $u$ to the line $l:\lambda u_1-u_2=0$. As $\lambda$ is irrational, the only integer point of $\sigma^\vee$ on line $l$ is $(0,0)$. Therefore for any point $u\neq 0$ in $S$, $\Delta(u)<\lambda$ and 
hence we can rewrite the distence function as
\begin{align*}
    D(u) = \frac{u_1(\lambda-\Delta(u))}{\sqrt{1+\lambda^2}}.
\end{align*}
Suppose $s_i = (s_i^1,s_i^2)\neq 0$ and for any $i<j$, we have $\Delta(s_i)>\Delta(s_j)$.  
If $i\le j$ then
\begin{align*}
    \Delta(s_i+s_j)-\Delta(s_i) = \frac{s_j^{1}(\Delta(s_j)-\Delta(s_i))}{s_i^{1}+s_j^1}\le 0
\end{align*} 
Then for an arbitrary non-zero element $u=\sum a_is_i\in S$, we obtain that
\begin{align*}
    D(u)\ge \frac{s^1_1(\lambda-\Delta(s_1))}{\sqrt{1+\lambda^2}} = D(s_1)>0.
\end{align*}
However, according to Kronecker's approximation theorem, $\{ \lambda u_1 - \lfloor \lambda u_1 \rfloor \mid u_1 \in \mathbb{Z}^+ \}$ is dense on 
$(0,1)$, hence we can choose a positive integer $u_1$ such that 
\begin{align*}
    0<\lambda u_1-\lfloor \lambda u_1 \rfloor<D(s_1)\sqrt{1+\lambda^2}
\end{align*}
and, consequently, $(u_1,\lfloor \lambda u_1 \rfloor)\in S$ and $D(u_1,\lfloor \lambda u_1 \rfloor)<D(s_1)$. The contradiction implies $S$ is not finitely generated.

Furthermore, since $S$ is not finitely generated, we can construct an infinite sequence of elements $\{{u_k}\} \subset S$ 
such that $D(u_{k+1}) < D(u_k)$. The corresponding chain of ideals $I_k = \langle \chi^{u_1}, \dots, \chi^{u_k} \rangle$ of semigroup 
algebra $\mathbb{C}[S]$ satisfies $I_1 \subsetneq I_2 \subsetneq I_3 \subsetneq \dots$, which never stabilizes. 
Therefore, $\mathbb{C}[S]$ is not Noetherian.
\end{solution}

\end{document}